\documentclass[10pt,a4paper]{amsart}
\usepackage[margin=19mm]{geometry}
\usepackage{amsmath,amssymb,mathtools}
\usepackage[scr=rsfs,cal=euler]{mathalfa}
\usepackage{xfrac}
\usepackage[unicode,hidelinks,bookmarks=false]{hyperref}
\newcommand{\ie}{i.e.\ }
\newcommand{\cf}{cf.\ }
\newcommand{\resp}{respectively\ }
\newcommand{\Subsection}{\subsection}
\providecommand{\nobreakdash}{\nobreak\hbox{-}}
\setlength{\emergencystretch}{2em}
\multlinegap0pt
\usepackage{bbm}
\usepackage{amssymb}
\usepackage{enumerate}
\usepackage[all]{xy}
\newtheorem{prop}{Proposition}
\newtheorem{cor}{Corollary}
\newtheorem{lem}{Lemma}
\newtheorem{defn}{Definition}
\newtheorem{thm}{Theorem}
\newtheorem{claim*}{Claim}
\title{K-moduli of quasimaps and on quasi-projectivity of K-moduli of Calabi-Yau fibrations over curves}
\author{Kenta Hashizume, Masafumi Hattori}
\date{Source: arXiv:2504.21519v2 (CC BY 4.0)}
\begin{document}
\maketitle

\noindent\emph{Source and licence.} K-moduli of quasimaps and on quasi-projectivity of K-moduli of Calabi-Yau fibrations over curves. Version arXiv:2504.21519v2 (CC BY 4.0), distributed by arXiv under the Creative Commons Attribution 4.0 International licence, http://creativecommons.org/licenses/by/4.0/, as recorded in the arXiv licence metadata for this exact version and shown on the abstract page https://arxiv.org/abs/2504.21519v2. Full licence notice: \texttt{licenses/arxiv-2504.21519-CC-BY-4.0.txt}.\par\smallskip

\noindent\textit{Editorial scope.} Counted unit: the article's thm thm--quasi--map--const (source0.tex lines 3877--3887). The article's complete original proof of it is delivered with it (source0.tex lines 3888--4015, one \texttt{proof} environment of 128 lines). No mathematics is written, repaired or re-derived: the statement and the proof are the article's own lines, in the article's own notation, and no retained line is substituted or re-flowed.  Retained uncounted as the source-local context the proof needs: lines 1267--1281; lines 1590--1593; lines 2201--2206; lines 2463--2466; lines 2801--2822; lines 3763--3777; lines 3849--3858; lines 4110--4116.  Nothing was omitted from any retained span, so the package carries no omission record.  The retained lines cite 2 works, whose entries are transcribed verbatim from the article's own bibliography source in the e-print and delivered with short editorial labels recorded in the item's reference mapping.  No retained line contains a figure environment, an image-inclusion command or drawing code, so no image is delivered and the package declares no asset dependency.  Editorial formatting only, in addition: an AMS \texttt{amsart} preamble that restates the article's own declarations and macros used by the retained lines, namely the environments \texttt{prop}, \texttt{cor}, \texttt{lem}, \texttt{defn}, \texttt{thm}, \texttt{claim*} on separate counters, together with the standard packages those lines need.  The retained environments print in this document as Proposition 1, Corollary 1, Lemma 1, Definition 1, Theorem 1, Claim* 1 in order of appearance, whereas the article prints the same items with the numbering of its own full document.  The complete native source of the article as distributed in the arXiv e-print for this exact version is bundled unchanged as \texttt{sources/177333/source0.tex}.
\begin{prop}\label{prop--smoothness--of--Abelian--moduli}
    Let $\mathcal{M}^{\mathrm{Ab}}_{d,v}\subset\mathcal{M}^{\mathrm{klt,CY}}_{d,v}$ be a substack such that 
     \[\mathcal{M}^{\mathrm{Ab}}_{d,v}(S)=
\left\{
 f\colon(\mathcal{X},\mathcal{L})\to S
\;\middle|
\begin{array}{l}
\text{$f\in\mathcal{M}^{\mathrm{klt,CY}}_{d,v}$ and for any geometric point $\bar{s}\in S$,}\\
\text{$\mathcal{X}_{\bar{s}}$ is an Abelian variety.}
\end{array}\right\}.
    \]
    Then, $\mathcal{M}_{d,v}^{\mathrm{Ab}}\subset\mathcal{M}^{\mathrm{klt,CY}}_{d,v}$ is smooth and an open and closed substack of $\mathcal{M}^{\mathrm{klt,CY}}_{d,v}$.

    In particular, the coarse moduli space $M_{d,v}^{\mathrm{Ab}}$ is normal and has only quotient singularities.
\end{prop}

\begin{cor}\label{cor--hk--moduli--normal}
    $\mathcal{M}^{\mathrm{symp}}_{2d,v}$ is smooth and an open and closed substack of $\mathcal{M}^{\mathrm{klt,CY}}_{2d,v}$.
    In particular, the coarse moduli space $M^{\mathrm{symp}}_{2d,v}$ of $\mathcal{M}^{\mathrm{symp}}_{2d,v}$ is normal with only quotient singularities.
\end{cor}

\begin{lem}\label{lem--stein--factorizaton}
Let $S$ be an arbitrary reduced scheme.
    For any object $f\colon (\mathcal{X},\mathscr{A})\to\mathcal{C}$ of $\mathcal{M}_{d,v,u,r}(S)$, the morphism $f$ is flat and we have $f_*\mathcal{O}_{\mathcal{X}}\cong\mathcal{O}_{\mathcal{C}}$.
    Furthermore, there exists a positive integer $l_0\in\mathbb{Z}_{>0}$, depending only on $d$, $v$, $u$, $r$, satisfying the following.
     There exists a line bundle $\mathscr{L}$ on $\mathcal{C}$ such that $f^*\mathscr{L}\cong \omega_{\mathcal{X}/S}^{[-l_0]}$.
  \end{lem}

\begin{lem}\label{lem--quasi--map--S_2}
    Let $X\hookrightarrow \mathbb{P}^N$ be an inclusion of a closed normal variety, let $\mathcal{C}\to S$ be a projective flat morphism to a normal variety with one-dimensional smooth geometric fibers, and let $q_1,q_2\colon\mathcal{C}\to[\mathrm{Cone}(X)\times S/\mathbb{G}_{m,S}]$ be two families of quasimaps such that $q_1$ is isomorphic to $q_2$ outside a closed subset $W\subset \mathcal{C}$ of codimension at least two.
    Then, $q_1$ and $q_2$ are globally isomorphic. 
\end{lem}

\begin{defn}\label{defn--associated--quasi--map--structure}
    Let $f\colon (W,A)\to\mathbb{P}_{\mathbbm{k}}^{1}$ be a uniformly adiabatically K-stable klt-trivial fibration over $\mathbb{P}_{\mathbbm{k}}^{1}$ of dimension $d$ with general fibers belonging to $\mathcal{M}_{d-1,v}^{\mathrm{klt},\mathrm{CY}}$.
    Let $U$ be a Zariski open subset of $\mathbb{P}_{\mathbbm{k}}^{1}$ such that $f|_{f^{-1}(U)}$ is smooth.
    Let $q^\circ\colon U\to M_{d-1,v}^{\mathrm{klt},\mathrm{CY}}$ be the moduli map corresponding to $f|_{f^{-1}(U)}$. 
    Let $B$ be the discriminant $\mathbb{Q}$-divisor and $M$ the moduli $\mathbb{Q}$-divisor associated to $f$.

    With notation as above, suppose that all the fibers of $f|_{f^{-1}U}$ are Abelian varieties or symplectic varieties smoothable to projective irreducible holomorphic symplectic varieties.
    Then, the image of $q^\circ$ is contained in a normal irreducible component $X^\circ$ of $\mathcal{M}_{d-1,v}^{\mathrm{klt},\mathrm{CY}}$ (Proposition \ref{prop--smoothness--of--Abelian--moduli} and Corollary \ref{cor--hk--moduli--normal}).
    Let $X$ be the Baily--Borel compactification of $X^\circ$ and $\Lambda_{\mathrm{Hodge}}$ the Hodge $\mathbb{Q}$-line bundle on $X$.
Take a sufficiently large and divisible positive integer $l$ such that $\Lambda_{\mathrm{Hodge}}^{\otimes l}$ is a very ample Cartier divisor and $lB$ is an integral divisor.    
Let $\iota\colon X\hookrightarrow \mathbb{P}_{\mathbbm{k}}^N$ be the closed embedding defined by $\Lambda_{\mathrm{Hodge}}^{\otimes l}$.
Then, we can extend $q^\circ$ to a quasimap 
$$q\colon \mathbb{P}_{\mathbbm{k}}^{1} \to [\mathrm{Cone}(X)/\mathbb{G}_{m, \mathbbm{k}}]$$
 of degree $l\cdot\mathrm{deg}(M+B)$ and of weight $\frac{1}{l}$ in the following way: 
Let $\mathrm{Cone}(X)\subset\mathbb{A}_{\mathbbm{k}}^{N+1}$ be the affine cone of $X$ as a subvariety of $\mathbb{P}_{\mathbbm{k}}^N$.
Let $\bar{q}\colon \mathbb{P}_{\mathbbm{k}}^1\to X$ be the unique extension of $q^{\circ}$ by the valuative criterion for properness.
By definition, we have $M\sim_{\mathbb{Q}}\bar{q}^*\Lambda_{\mathrm{Hodge}}$. 
Now take the canonical coordinate sections $x_0,\ldots,x_N$ of $\mathbb{P}^N_{\mathbbm{k}}$ and we set $q$ as the quasimap defined by the sections $\bar{q}^*s_i\cdot s_{lB}\in H^0(\mathbb{P}_{\mathbbm{k}}^{1},\mathcal{O}_{\mathbb{P}_{\mathbbm{k}}^{1}}(l(M+B)))$, where $s_{lB}\in H^0(\mathbb{P}_{\mathbbm{k}}^{1},\mathcal{O}_{\mathbb{P}_{\mathbbm{k}}^{1}}(lB))$ is the section corresponding to $lB$. 

We call the $q\colon \mathbb{P}_{\mathbbm{k}}^{1} \to [\mathrm{Cone}(X)/\mathbb{G}_{m, \mathbbm{k}}]$ a {\it quasimap associated with} $f$.
We note that $q$ depends on the integer $l$, and if we fix $l$ then such $q$ is uniquely determined.
\end{defn}

\begin{thm}\label{thm--Fujino's--period--mapping--theory}
    We put $l=gkl_0$, where $l_0$ as in Lemma \ref{lem--stein--factorizaton}. 
    Note that $l$ is a positive integer depending only on $d$, $v$, $u$, $r$, and $V$. 
    Then the the following holds.
    Let $S$ be a smooth affine variety and $f\colon (\mathcal{X},\mathcal{A})\to\mathcal{P}$ a family belonging to $\mathcal{M}_{d,v,u,r,V}^{\mathrm{CYFib}}(S)$ such that $\mathcal{A}$ is represented by a line bundle on $\mathcal{X}$.
    Let $\mathcal{P}_{\mathrm{klt}}$ be the largest open subset of $\mathcal{P}$ such that all geometric fibers of $f$ over $\mathcal{P}_{\mathrm{klt}}$ are klt. 
    Let $\mu^\circ\colon \mathcal{P}_{\mathrm{klt}}\to V$ be the morphism induced by the family $f|_{f^{-1}(\mathcal{P}_{\mathrm{klt}})}$ of klt Calabi--Yau varieties. 
    Take an Ambro model $\rho\colon\widetilde{\mathcal{P}}\to \mathcal{P}$ with respect to $f$, and let $\widetilde{\mathcal{M}}$ and $\widetilde{\mathcal{B}}$ be the moduli $\mathbb{Q}$-divisor and the discriminant $\mathbb{Q}$-divisor on $\mathcal{P}$, respectively.
    Then, after replacing $\widetilde{\mathcal{P}}$ with a higher birational model suitably, we have:
    \begin{enumerate}
       \item \label{thm--Fujino's--period--mapping--theory-(1)} $l\widetilde{\mathcal{M}}$ and $l\widetilde{\mathcal{B}}$ are both $\mathbb{Z}$-divisors, and
        \item \label{thm--Fujino's--period--mapping--theory-(2)} putting $L:=h^0(\overline{V}^{\mathrm{BB}},\Lambda^{\otimes l}_{\mathrm{Hodge}})-1$, then there exist $L+1$ sections $$\widetilde{\varphi}_0,\ldots,\widetilde{\varphi}_L\in H^0(\widetilde{\mathcal{P}},\mathcal{O}_{\widetilde{\mathcal{P}}}(l\widetilde{\mathcal{M}}))$$ that generate $l\widetilde{\mathcal{M}}$ and define a morphism $\widetilde{\mu}\colon\widetilde{\mathcal{P}}\to \mathbb{P}^L$ such that $\widetilde{\mu}^*\mathcal{O}(1)\sim l\widetilde{\mathcal{M}}$, $\widetilde{\mu}$ factors $\iota\colon \overline{V}^{\mathrm{BB}}\hookrightarrow \mathbb{P}^L$ and this is equivalent to the composition of  $\rho$, $\mu^\circ$ and $\iota$ as rational maps.  
        Here, $\iota\colon \overline{V}^{\mathrm{BB}}\to \mathbb{P}^L$ is induced by the complete linear system $|\Lambda^{\otimes l}_{\mathrm{Hodge}}|$ $=|\iota^*\mathcal{O}_{\mathbb{P}^L}(d'l_{0})|$. 
    \end{enumerate}
\end{thm}

\begin{lem}\label{lemma--description--of--base--locus}
     Let $S$ be a smooth quasi-projective variety, and let $f\colon (\mathcal{X},\mathcal{A})\to\mathcal{P}$ be a family that belongs to $\mathcal{M}_{d,v,u,r,V}^{\mathrm{CYFib}}(S)$ such that $\mathcal{A}$ is represented by a line bundle.
    Take an arbitrary Ambro model $\rho\colon\tilde{\mathcal{P}}\to \mathcal{P}$ with respect to $f$, and let $\widetilde{\mathcal{M}}$ and $\widetilde{\mathcal{B}}$ be the moduli $\mathbb{Q}$-divisor and the discriminant $\mathbb{Q}$-divisor, respectively.
    Let $E$ be the $\rho$-exceptional divisor $\mathbb{Q}$-linearly equivalent to $K_{\widetilde{\mathcal{P}}}-\rho^*K_{\mathcal{P}}$. 
    Then, 
    \[
E+\widetilde{\mathcal{B}}+\widetilde{\mathcal{M}}=\rho^*(\mathcal{B}+\mathcal{M})
\]
    holds and $\widetilde{\mathcal{B}}+E$ is effective. 
\end{lem}

\begin{defn}\label{defn--isotrivial}
    Let $f\colon X\to S$ be a projective morphism of schemes of finite type over $\mathbb{C}$.
    We say that $X$ is {\it isotrivial} over $S$ if $X_s$ and $X_{s'}$ are isomorphic for any closed points $s,s'\in S$.

    Let $q\colon (C,B)\to [\mathbb{A}^{N+1}_S/\mathbb{G}_{m,S}]$ be a family of log Fano quasimaps over $S$.
    We say that $q$ is {\it isotrivial} over $S$ if $q_s$ and $q_{s'}$ are isomorphic for any closed points $s,s'\in S$.
\end{defn}

\begin{thm}\label{thm--quasi--map--const}
        Let $S$ be a quasi-projective normal variety and let $f\colon(\mathcal{X},\mathcal{A})\to\mathcal{P}$ be an object of $\mathcal{M}_{d,v,u,r,V}^{\mathrm{CYFib}}(S)$.
        Suppose that $\mathcal{A}$ is a line bundle.
Then, there exists the following object $q\colon \mathcal{P}\to [\mathrm{Cone}(\overline{V}^{\mathrm{BB}})\times S/\mathbb{G}_{m,S}]$ of $\mathcal{M}_{l(2-u),1,\frac{1}{l},u,\iota}^{\mathrm{Kst.qmaps}}(S)$ uniquely up to isomorphism satisfying the following. 
\begin{enumerate}
\item Let $\mathscr{L}$ be the associated line bundle on $\mathcal{P}$ with $q$. Then, $f^*\mathscr{L}\sim_S l(K_{\mathcal{X}/S}-f^*K_{\mathcal{P}/S})$.
\item For any geometric point $\bar{s}\in S$, let $B_{\bar{s}}$ be the fixed part of $q_{\bar{s}}$.
Then, $\frac{1}{l}B_{\bar{s}}$ is the discriminant $\mathbb{Q}$-divisor with respect to $f_{\bar{s}}$ and $q_{\bar{s}}$ defines the moduli map associated to $f_{\bar{s}}$ (cf.~Definition \ref{defn--associated--quasi--map--structure}).
\end{enumerate}
Furthermore, $f^*\mathscr{L}\sim l(K_{\mathcal{X}/S}-f^*K_{\mathcal{P}/S})$.
    \end{thm}

    \begin{proof}
        We first deal with the uniqueness of $q$.
        Assume that there exists another family of quasimaps $q'$ with the same property as $q$.
        Then, $q'$ and $q$ define two line bundles $\mathscr{L}'$ and $\mathscr{L}$.
        By the assumption, we see that $\mathscr{L}_{\bar{s}}\sim\mathscr{L}'_{\bar{s}}$ for any geometric point $\overline{s}\in S$.
Since it is enough to show that $q$ and $q'$ are isomorphic to each other Zariski locally on $S$ and all geometric fibers of $\mathcal{P}\to S$ are irreducible and smooth, we may assume that $\mathscr{L}\sim\mathscr{L}'$ by shrinking $S$.
Let $\varphi_0,\ldots,\varphi_L$ (resp.~$\varphi'_0,\ldots,\varphi'_L$) be the sections corresponding to the canonical coordinates of $\mathbb{A}^{L+1}$ of $\mathscr{L}$ (resp.~$\mathscr{L}'$).
By the assumption, we see that for any geometric point $\bar{s}\in S$, $\varphi_{0,\bar{s}},\ldots,\varphi_{L,\bar{s}}$ and $\varphi'_{0,\bar{s}},\ldots,\varphi'_{L,\bar{s}}$ define the same quasimap structures.
This is equivalent to that there exists $\alpha_{\bar{s}}\in\overline{\kappa(s)}^{\times}$ such that $\alpha_{\bar{s}}=\frac{\varphi_{j,\bar{s}}}{\varphi'_{j,\bar{s}}}$ for any $j=0,\ldots,L$ such that $\varphi_{j,\bar{s}}\cdot \varphi'_{j,\bar{s}}\ne0$.
This means that shrinking $S$ if necessary, we may assume that there exists $\alpha\in H^0(S,\mathcal{O}_S^{\times})$ such that $\varphi_i=\alpha \varphi'_i$ for any $i\in\{0,\ldots,L\}$.
Hence, $q$ and $q'$ are isomorphic families of quasimaps to each other.

Next, we deal with the existence of such $q$ as in the assertion.
To show this, we first show that we may replace $S$ with any log resolution of $S$ freely.
More precisely, we show the following claim.
\begin{claim*}
    Let $\varphi\colon S'\to S$ be an arbitrary projective birational morphism from a smooth variety.
    Suppose that there exists a family of quasimaps $q'\colon \mathcal{P}_{S'}\to [\mathrm{Cone}(\overline{V}^{\mathrm{BB}})\times S'/\mathbb{G}_{m,S'}]$ as the assertion for $f_{S'}\colon (\mathcal{X}\times_SS',\mathcal{A}_{S'})\to \mathcal{P}_{S'}$.
    Then, there exists a family of quasimaps $q$ as the assertion for $S$ such that $q_{S'}$ is isomorphic to $q'$.
\end{claim*}

\begin{proof}[Proof of Claim]   
Let $\mathscr{L}'$ be the associated line bundle on $\mathcal{P}_{S'}$ with $q'$ and $\varphi'_0,\ldots,\varphi'_{L}\in\mathrm{Sec}_{q'}(\mathscr{L}')$ the sections corresponding to the canonical basis.
Note that by uniqueness of $q$ as the assertion, it suffices to show the existence of $q$ locally over $S$.
    By the property of $q'$, we see that $f_{S'}^*\mathscr{L}'\sim_{S'}l(K_{\mathcal{X}_{S'}/S'}-f_{S'}^*K_{\mathcal{P}_{S'}/S'})$.
   Since $f\in \mathcal{M}_{d,v,u,r}(S)$, $\omega_{\mathcal{X}/S}^{[-l]}$ is a pullback of a certain relatively ample line bundle on $\mathcal{P}$ over $S$. 
    By the choice of $l$ and Lemma \ref{lem--stein--factorizaton}, there exists a line bundle $\mathscr{L}$ on $\mathcal{P}$ such that $f^*\mathscr{L}=l(K_{\mathcal{X}/S}-f^*K_{\mathcal{P}/S})$. 
    By Lemma \ref{lem--stein--factorizaton} and the assumption that $S'$ is smooth, we know that $f_{S'*}\mathcal{O}_{\mathcal{X}_{S'}}\cong \mathcal{O}_{\mathcal{P}_{S'}}$.
    We also have that $\varphi_{\mathcal{P}*}\mathcal{O}_{\mathcal{P}_{S'}}\cong\mathcal{O}_{\mathcal{P}}$ for $\varphi_{\mathcal{P}}\colon\mathcal{P}_{S'}\to \mathcal{P}$. 
    Thus, $(\varphi_{\mathcal{P}}\circ f_{S'})_*\mathcal{O}_{\mathcal{X}_{S'}}\cong \mathcal{O}_{\mathcal{P}}$ holds and $\mathscr{L}'\sim_{S'} \varphi_{\mathcal{P}}^*\mathscr{L}$ by \cite[Lemma 9.2.7]{FGA}.
    Here, we show that $\mathscr{L}'\sim_S\varphi_{\mathcal{P}}^*\mathscr{L}$.
    Indeed, let $\mathscr{M}$ be a line bundle on $S'$ such that $\mu'^*\mathscr{M}\sim\mathscr{L}'\otimes\varphi_{\mathcal{P}}^*\mathscr{L}^{\otimes -1}$, where $\mu'\colon \mathcal{P}_{S'}\to S'$ is the canonical morphism.
    Take an arbitrary closed point $s\in S$ and an arbitrary closed point $p\in \mathcal{P}_s$.
    Let $U\subset \mathcal{P}$ be an open affine neighborhood of $p$ such that $\mathscr{L}|_U\sim\mathcal{O}_U$ and $\mu'^*\mathscr{M}|_{\varphi_{\mathcal{P}}^{-1}(U)}\sim\mathscr{L}'|_{\varphi_{\mathcal{P}}^{-1}(U)}$.
    By shrinking $U$ and renumbering $i$ if necessary, we may assume that for any $s_1\in U$, there exists a point $p_1\in \varphi_{\mathcal{P}}^{-1}(s_1)$ such that $\varphi'_{0,p_1}\ne0$.
    Here, by the definition of $q'$, the family $q'|_{\varphi^{-1}(\varphi(\mu'(s_1)))}$ is isotrivial (see Definition \ref{defn--isotrivial} below).
    Therefore, it is easy to see that $\varphi'_0$ does not vanish at any point of $\varphi_{\mathcal{P}}^{-1}(U)$.
    This means that $\mu'^*\mathscr{M}\sim_{\mathcal{P}}0$.
    By this, it is not hard to see that $\mathscr{M}'\sim_S0$.
   Thus, we may assume that $\mathscr{L}'\sim\varphi_{\mathcal{P}}^*\mathscr{L}$ by shrinking $S$. 
    Via the identification $H^0(\mathcal{P}',\mathscr{L}')\cong H^0(\mathcal{P},\mathscr{L})$, we can consider $\varphi'_i$'s to be the sections in $H^0(\mathcal{P},\mathscr{L})$.
    Then, $\varphi'_i\in H^0(\mathcal{P},\mathscr{L})$ define a family of quasimaps $q$ such that $q_{S'}=q'$.
    By this property, it is easy to check that $q$ has the desired properties.
    We complete the proof. 
\end{proof}

\noindent{\it Proof of Theorem \ref{thm--quasi--map--const} continued.} 
Thus, replacing $S$ with a suitable log resolution $S'$ of $S$ if necessary, we may assume that there exist $\rho\colon\tilde{\mathcal{P}}\to\mathcal{P}$ and $L+1$ sections $\widetilde{\varphi_0},\ldots,\widetilde{\varphi_L}\in H^0(\tilde{\mathcal{P}},\mathcal{O}_{\widetilde{\mathcal{P}}}(l\widetilde{\mathcal{M}}))$ as in Theorem \ref{thm--Fujino's--period--mapping--theory}.
Here, $\widetilde{\mathcal{M}}$ and $\widetilde{\mathcal{B}}$ are the moduli $\mathbb{Q}$-divisor and the discriminant $\mathbb{Q}$-divisor on $\tilde{\mathcal{P}}$ respectively.
Let $E$ be the $\rho$-exceptional divisor $\mathbb{Q}$-linearly equivalent to $K_{\widetilde{\mathcal{P}}}-\rho^*K_{\mathcal{P}}$. 
By the choice of $l$ and Lemma \ref{lem--stein--factorizaton}, there exists a line bundle $\mathscr{L}$ on $\mathcal{P}$ such that $f^*\mathscr{L}=l(K_{\mathcal{X}/S}-f^*K_{\mathcal{P}/S})$.
In this section, we construct the desired family of quasimaps by using the $L+1$ sections.
By Lemma \ref{lemma--description--of--base--locus}, there exists the natural injection
\begin{equation}\tag{$\clubsuit$}\label{eq--injection-of--function}
H^0(\tilde{\mathcal{P}},\mathcal{O}_{\widetilde{\mathcal{P}}}(l\widetilde{\mathcal{M}}))\hookrightarrow H^0(\tilde{\mathcal{P}},\mathcal{O}_{\widetilde{\mathcal{P}}}(l(\widetilde{\mathcal{M}}+\widetilde{\mathcal{B}}+E)))\cong H^0(\mathcal{P},\mathscr{L}).
\end{equation}
Via this identification, we identify $\widetilde{\varphi_i}$ with a section $\varphi_i\in H^0(\mathcal{P},\mathscr{L})$ for any $i$.
We claim that $\varphi_i\in H^0(\mathcal{P},\mathscr{L})$ define a structure of family of quasimaps $q\colon \mathcal{P}\to[\mathrm{Cone}(\overline{V}^{\mathrm{BB}})\times S/\mathbb{G}_{m,S}]$.
Indeed, let $\mathcal{P}_{\mathrm{klt}}$ be the open subset of $\mathcal{P}$ such that any geometric fiber of $f$ over $\mathcal{P}_{\mathrm{klt}}$ is klt and $\mu^\circ\colon \mathcal{P}_{\mathrm{klt}}\to \overline{V}^{\mathrm{BB}}$ the morphism induced by the family $f|_{f^{-1}(\mathcal{P}_{\mathrm{klt}})}$ of klt Calabi--Yau varieties.
Let $D$ be a prime divisor on $S$.
We note that the canonical morphism $\mathcal{P}_{\mathrm{klt}}\to S$ is surjective and at any general point of $s\in D$, some $\varphi_{i}|_{\mathcal{P}_{\mathrm{klt},s}}$ does not vanish since $\mathcal{P}_{\mathrm{klt},s}$ is not contained in $\rho(\mathrm{Exc}(\rho))$.
This shows that $q$ is a family of quasimaps over $S$ in codimension one.
For any closed point $s\in S$, take the one point blow-up $\xi\colon \mathrm{Bl}_sS\to S$.
As Theorem \ref{thm--Fujino's--period--mapping--theory}, there exist $\rho'\colon \widetilde{\mathcal{P}}'\to\mathcal{P}_{\mathrm{Bl}_sS}$ and $L+1$ sections $\widetilde{\varphi_0}',\ldots,\widetilde{\varphi_L}'\in H^0(\tilde{\mathcal{P}}',\mathcal{O}_{\widetilde{\mathcal{P}}'}(l\widetilde{\mathcal{M}}'))$ corresponding to the canonical coordinate of $\mathbb{A}^{L+1}$.
Here, $\widetilde{\mathcal{M}}'$ and $\widetilde{\mathcal{B}}'$ are the moduli $\mathbb{Q}$-divisor and the discriminant $\mathbb{Q}$-divisor on $\tilde{\mathcal{P}}'$ respectively.
Let $E'$ be the $\rho'$-exceptional divisor $\mathbb{Q}$-linearly equivalent to $K_{\widetilde{\mathcal{P}}'}-\rho'^*K_{\mathcal{P}_{\mathrm{Bl}_sS}}$. 
Let $\widetilde{\mathcal{W}}$ be a smooth variety such that there exist birational projective morphisms $p_{\tilde{\mathcal{P}}}\colon \widetilde{\mathcal{W}}\to \tilde{\mathcal{P}}$ and $p_{\widetilde{\mathcal{P}}'}\colon \widetilde{\mathcal{W}}\to \widetilde{\mathcal{P}}'$.
Note that $p_{\widetilde{\mathcal{P}}'}^*\widetilde{\mathcal{M}}'\sim p_{\tilde{\mathcal{P}}}^*\widetilde{\mathcal{M}}$ by the definition of Ambro models.
Therefore, we can show that the natural injection
\begin{align*}
   H^0(\widetilde{\mathcal{P}}',\mathcal{O}_{\widetilde{\mathcal{P}}'}(l\widetilde{\mathcal{M}}'))\hookrightarrow H^0(\widetilde{\mathcal{P}}',\mathcal{O}_{\widetilde{\mathcal{P}}'}(l(\widetilde{\mathcal{M}}'+\widetilde{\mathcal{B}}'+E')))\cong H^0(\mathcal{P},\mathscr{L})
\end{align*}
obtained as \eqref{eq--injection-of--function} is naturally isomorphic to \eqref{eq--injection-of--function} and by this identification, $\widetilde{\varphi_i}'$ coincides with $\widetilde{\varphi_i}$ for any $i$.
By this fact, since $\widetilde{\varphi_i}'$ does not vanish over the generic point of $\xi^{-1}(s)$ for some $i$ as we have shown, $\widetilde{\varphi_i}_s$ does not vanish.
Thus, $q$ is a family of quasimaps over $S$.
Note that if $q$ satisfies conditions (1) and (2), then $f^*\mathscr{L}\sim l(K_{\mathcal{X}/S}-f^*K_{\mathcal{P}/S})$ holds by construction.

 It suffices to show that $q$ satisfies the condition (2) for any geometric point $\bar{s}\in S$.
 In this section, we reduce this to the case when $S$ is a smooth curve and the condition (2) holds for all but one closed points of $S$.
We first note that it is enough to check condition (2) only for closed points for $S$.
Indeed, it is easy to check that for any subvariety $W\subset S$, if an arbitrary closed point of $W$ satisfies condition (2), then condition (2) also holds for the geometric generic point of $W$ using the same argument as the proof of \cite[Lemma 4.1]{HH}. 
Indeed, by the proof of \cite[Lemma 4.1]{HH}, we see that there exists an \'etale morphism $W'\to W$ such that for  $\mathcal{P}_{W'}$, we can take an Ambro model $\rho\colon\widetilde{\mathcal{P}}\to \mathcal{P}_{W'}$ such that $(\widetilde{\mathcal{P}},\widetilde{\mathcal{M}}+\widetilde{\mathcal{B}}+\mathrm{Exc}(\rho))$ is relatively simple normal crossing over $W'$, where $\widetilde{\mathcal{M}}$ and $\widetilde{\mathcal{B}}$ are the moduli $\mathbb{Q}$-divisor and the discriminant $\mathbb{Q}$-divisor respectively.
We also note that for any general closed point $s\in S$, we can easily check that condition (2) holds using the same argument as the proof of \cite[Lemma 4.1]{HH}.
For any closed point $s\in S$,  we take a morphism $g\colon C\to S$ from a smooth affine curve such that a general point of $C$ is mapped to a general point of $S$ and a closed point $0\in C$ is mapped to $s$ under $g$.
Consider $g^*q$ and $q_C$.
Here, the latter is the quasimap constructed for $C$ in the same way as $q$ for $S$.
We claim that $g^*q$ and $q_C$ are isomorphic.
Indeed, there exists a non-empty open subset $U\subset C$ such that $g^*q$ and $q_C$ are isomorphic over $\mathcal{P}_{\mathrm{klt}}\times_SC\cup \mathcal{P}\times_SU$.
Note that $$\mathrm{codim}_{\mathcal{P}\times_SC}(\mathcal{P}\times_SC\setminus(\mathcal{P}_{\mathrm{klt}}\times_SC\cup \mathcal{P}\times_SU))\ge2.$$ 
Thus, $g^*q$ and $q_C$ are canonically isomorphic to each other by Lemma \ref{lem--quasi--map--S_2}.
Therefore, replacing $S$ with $C$ and shrinking $C$ if necessary, we may assume that $C=S$ is a curve with a closed point $0\in C$ such that at least the condition (2) holds for any closed point of $C\setminus\{0\}$.

Finally, we deal with that under the assumption that that $S$ is a curve with a closed point $0\in S$ such that at least the condition (2) holds for any closed point of $S\setminus\{0\}$, 
$\varphi_0|_0,\ldots,\varphi_L|_0$ define the canonical quasimap structure associated with $(\mathcal{X}_0,\mathcal{A}_0)\to \mathbb{P}^1$ (cf.~Definition \ref{defn--associated--quasi--map--structure}).
Since $\mathcal{P}$ is a ruled surface over $S$, we may assume that $\mathcal{P}\cong\mathbb{P}^1\times S$ by shrinking $S$. 
 Take an Ambro model $\rho\colon\widetilde{\mathcal{P}}\to \mathbb{P}^1\times S$ of $\mathcal{X}\to\mathbb{P}^1\times S$, where $\rho\colon\widetilde{\mathcal{P}}\to \mathbb{P}^1\times S$ is a projective birational morphism from a smooth variety as in Theorem \ref{thm--Fujino's--period--mapping--theory}.
    Let $D\cong\mathbb{P}^1$ be the strict transform of $\mathbb{P}^1\times\{0\}$ in $\mathcal{P}$. 
    Note that $D$ is isomorphic to $\mathbb{P}^1\times\{0\}$ via $\rho$.
    Let $(\widetilde{\mathcal{X}},\widetilde{\mathcal{D}})\to\widetilde{\mathcal{P}}$ be a projective contraction of smooth varieties that is birational to $\mathcal{X}\to \mathbb{P}^1\times S$ such that $\mathcal{X}$ and $(\widetilde{\mathcal{X}},\widetilde{\mathcal{D}})$ are log crepant.
    Let $\widetilde{\mathcal{B}}$ and $\widetilde{\mathcal{M}}$ be the discriminant and the moduli $\mathbb{Q}$-divisors associated with $(\widetilde{\mathcal{X}},\widetilde{\mathcal{D}})\to\widetilde{\mathcal{P}}$ respectively. 
    We note that $E:=K_{\widetilde{\mathcal{P}}}-\rho^*K_{\mathbb{P}^1\times S}$ is $\rho$-exceptional and effective since $S\times \mathbb{P}^1$ is smooth. 
    By replacing $\rho$ with a further blow up, we may assume that $(\widetilde{\mathcal{P}},\widetilde{\mathcal{B}}+\widetilde{\mathcal{M}}+\widetilde{\mathcal{P}}_0 + \mathrm{Exc}(\rho))$ is simple normal crossing.
    Here, the sections $\varphi_0,\ldots,\varphi_L$ define the morphism $\mu^\circ\colon (\mathbb{P}^1\times S)_{\mathrm{klt}}\to V$ by the construction of $q$.
    Furthermore, if $\tilde{\mu}\colon \widetilde{\mathcal{P}}\to \overline{V}^{\mathrm{BB}}$ denotes the morphism that is equivalent to $\mu^\circ$ as a rational map, then via the identification 
    \[ H^0(\tilde{\mathcal{P}},\mathcal{O}_{\widetilde{\mathcal{P}}}(l\widetilde{\mathcal{M}}))\ni\rho^*\varphi\hookrightarrow  \varphi\in H^0(\mathcal{P},\mathscr{L}),
    \]
    $\rho^*\varphi_i$'s generate $l\widetilde{\mathcal{M}}$ and define $\tilde{\mu}$.
    Thus, if we regard $\rho^*\varphi_i$ as a section of $H^0(\widetilde{\mathcal{P}},\rho^*\mathscr{L})$, then $\rho^*\varphi_i$'s generate the ideal sheaf $\mathcal{O}_{\widetilde{\mathcal{P}}}(-l(\widetilde{\mathcal{B}}+E))$ (cf.~Lemma \ref{lemma--description--of--base--locus}).
    Let $(\widehat{\mathcal{X}},\widehat{\mathcal{D}})$ be the main component of $(\widetilde{\mathcal{X}},\widetilde{\mathcal{D}}+\widetilde{\mathcal{X}}_0-\widehat{\mathcal{X}})\times_{\widetilde{\mathcal{P}}}D$.
    Here, we may assume that $\widehat{\mathcal{X}}$ is a smooth divisor in $\widetilde{\mathcal{X}}$.
    By restricting $\tilde{\mu}$ to $D$, we have that $\widetilde{\mathcal{M}}|_D$ is the moduli $\mathbb{Q}$-divisor associated with the subklt--trivial fibration $(\widehat{\mathcal{X}},\widehat{\mathcal{D}})\to D$. 
    It is easy to see that $(\widehat{\mathcal{X}},\widehat{\mathcal{D}}+(\widetilde{\mathcal{X}_0}-\widehat{\mathcal{X}})|_{\widehat{\mathcal{X}}})$ is log crepant to $\mathcal{X}_0$ by the adjunction formula.
    Thus, $(D,(\widetilde{\mathcal{M}}+\widetilde{\mathcal{B}}+\widetilde{\mathcal{P}}_0-D)|_D)$ is isomorphic to the generalized log pair defined by $\mathcal{X}_0\to\mathbb{P}^1$.
    Since $\widetilde{\mathcal{M}}|_D$ is the moduli $\mathbb{Q}$-divisor, it is not hard to see that $(\widetilde{\mathcal{B}}+\widetilde{\mathcal{P}}_0-D)|_D$ is the discriminant divisor with respect to $\mathcal{X}_0\to\mathbb{P}^1$.
    Furthermore, we note that $E|_D=(\widetilde{\mathcal{P}}_0-D)|_D$ as Cartier divisors.
    Indeed, consider the following equation of log canonical divisors
    \[
    K_{\widetilde{\mathcal{P}}}+\widetilde{\mathcal{P}}_{0}=\rho^*(K_{\mathbb{P}^1\times S}+\mathbb{P}^1\times\{0\})+E.
    \]
    By this, we have
    \[
    K_{D}+(\widetilde{\mathcal{P}}_{0}-D)|_D=\rho|_D^*K_{\mathbb{P}^1}+E|_D.
    \]
    Hence, $E|_D=(\widetilde{\mathcal{P}}_0-D)|_D$ as Cartier divisors.
    Since $\rho^*\varphi_0|_D,\ldots,\rho^*\varphi_L|_D$ generate the ideal sheaf $\mathcal{O}_{D}(-l(\widetilde{\mathcal{B}}+E)|_D)$, we see that $\varphi_0|_0=\rho^*\varphi_0|_D,\ldots,\varphi_L|_0=\rho^*\varphi_L|_D$ define the associated quasimap structure to the klt--trivial fibration $f_0\colon \mathcal{X}_0\to\mathbb{P}^1$.
    This shows that condition (2) holds in this case.
    We complete the proof.
    \end{proof}

\begin{thebibliography}{99}
\bibitem[FGA]{FGA} B. Fantechi, et. al. {\it Fundamental Algebraic Geometry: Grothendieck's FGA Explained}. Amer. Math. Soc. 2005.
\bibitem[HH]{HH} K. Hashizume, M. Hattori, On boundedness and moduli spaces of K-stable Calabi-Yau fibrations over curves, Geom.~Topol. (3), {\bf29} (2025), 1619--1691.
\end{thebibliography}
\end{document}
